% Conclusions and Recommendations (BGU \S 10.7.2.9).
% V5: restores the BGU-mandated project-management content that the V4
% peer-review revision had stripped --- an explicit scientific conclusion,
% a goal-by-goal evaluation against the preliminary (PDR) targets with the
% architecture trade-offs, and a plan-versus-execution comparison (schedule
% and compute budget). The continuation recommendations are retained.
\section{Conclusions and Recommendations}
\label{sec:conclusions}

\subsection{Scientific conclusions}
This work delivered a reproducible 402-cell training campaign across three
architectures (ResNet50, DenseNet121, TransNeXt-tiny), two datasets (CIFAR-10,
MNIST), five degradation levels, and six experimental phases, anchored to a
deterministic per-pipeline PSNR / SSIM probe. Four conclusions follow from the
measurements.

\begin{enumerate}
  \item \textbf{The collapse is an information bottleneck, not over-fitting.}
        Layering dropout, mix-up and cut-mix on the combined-degradation
        pipeline (Phase~D) recovers only a small fraction of the lost accuracy,
        and simplifying the protocol to the THz-relevant operators (Phase~B2)
        recovers a larger but still modest margin. Because training-side
        intervention cannot close the gap, the dominant cause is destroyed
        information rather than a generalization failure.
  \item \textbf{The foveal transformer is the most robust backbone at the
        operationally relevant degradation levels.} TransNeXt-tiny holds the
        largest accuracy margin over the two CNNs at the mild and moderate end
        of the combined curve, and is the most robust architecture on every
        single-axis comparison at matched image quality.
  \item \textbf{Architectures converge at extreme degradation.} At the Phase~B1
        L5 operating point on CIFAR-10 all three backbones fall into a tight
        $[19,21]\,\%$ band: once the spatial information is removed, the
        architectural choice ceases to matter.
  \item \textbf{Resolution is the dominant degradation axis.} On the PSNR and
        SSIM scales the resolution axis carries by far the strongest
        accuracy--quality association (Spearman $\rho = +0.98$); blur and
        salt-and-pepper are markedly weaker. Restoring spatial resolution is
        therefore the highest-leverage front-end intervention.
\end{enumerate}

\subsection{Project evaluation against the preliminary goals}
The preliminary design report (PDR) set three acceptance criteria. The project
is evaluated against each below.

\begin{itemize}
  \item \textbf{Goal 1 --- at least $80\%$ top-1 accuracy at the moderate (L3)
        operating point.} \emph{Partially met, and informatively so.} On the
        full combined L3 pipeline the target is not reached on CIFAR-10
        (accuracies cluster near $36$--$39\%$); it is essentially met on MNIST
        (about $79$--$82\%$) and is comfortably exceeded for the milder
        single-axis operating points. The shortfall on combined CIFAR-10 L3 is
        not an engineering failure but the central scientific result: the
        moderate combined degradation destroys enough information that no
        backbone in the design space reaches $80\%$.
  \item \textbf{Goal 2 --- an identifiable architectural winner beating the
        baseline by a meaningful margin.} \emph{Met in the relevant regime.}
        TransNeXt-tiny is the clear winner at mild and moderate degradation,
        and the margin exceeds the multi-seed noise floor ($\sigma<1.5$~\Mpp{}).
        The margin vanishes at L5, so the ``winner'' claim is scoped to the
        mild-to-moderate range where it is decision-relevant.
  \item \textbf{Goal 3 --- inference latency compatible with sensor-side
        deployment.} \emph{Met for the CNNs.} ResNet50 and DenseNet121 are
        comfortably within the latency budget; TransNeXt-tiny is roughly
        $1.5\times$ slower because its attention path is not graph-compiled on
        the current stack, which is the subject of a continuation recommendation
        below.
\end{itemize}

\textbf{Trade-offs of the final architecture choice.}\quad ResNet50 is the
fastest and simplest backbone and an excellent robust baseline, but it carries
no robustness advantage; DenseNet121 gains a little at moderate CIFAR-10
degradation through feature reuse, at the cost of higher memory and greater
chroma-noise sensitivity; TransNeXt-tiny delivers the best accuracy at mild and
moderate degradation through foveal attention, at the cost of higher latency and
parameter count and with no advantage at extreme degradation. For a THz
deployment operating at mild-to-moderate degradation, TransNeXt-tiny is the
recommended backbone; at extreme degradation the engineering effort is better
spent on front-end restoration than on the classifier.

\subsection{Plan versus execution}
The project was delivered within its planned schedule and compute envelope.
Table~\ref{tab:plan_exec} compares the PDR / PRD compute budget with the
campaign's logged execution. The single material overrun was Phase~D
($+7.6\%$ over budget, attributable to the early-stopping patience triggering
later than estimated on several cells); it was absorbed by the
$\sim\!310$~GPU-hour saving from substituting \texttt{transnext\_tiny} for the
originally planned \texttt{transnext\_small}, so the campaign finished well
inside the original envelope. All scheduled milestones --- PDR, software and
determinism-gate implementation, the 402-cell campaign, the multi-seed audit and
post-training diagnostics, and this final report --- were met against the
26~July~2026 submission deadline.

\begin{table}[htbp]
\centering
\caption{Plan versus execution: compute budget. Planned figures are the PDR / PRD estimates; actual figures are the logged GPU-hours recorded in \texttt{Final\_Exp.json}.}
\label{tab:plan_exec}
\setlength{\tabcolsep}{5pt}
\renewcommand{\arraystretch}{1.15}
\begin{tabular}{lrrl}
\toprule
Work item & Planned (GPU-h) & Actual (GPU-h) & Status \\
\midrule
Hyperparameter pre-tuning (Optuna)        & 20.0  & 19.4  & Within budget \\
Phases A / B1 / C (186 cells)             & 62.0  & 61.2  & Within budget \\
Phase D regularization sweep (90 cells)   & 52.5  & 56.5  & $+7.6\%$ over \\
Phases B2 / B2-nr / C2 (126 cells)        & 66.5  & 65.8  & Within budget \\
Multi-seed audit (48 replicates)          & 24.0  & 23.6  & Within budget \\
Post-training diagnostics (no re-train)   & 3.5   & 3.4   & Within budget \\
\midrule
Total                                     & 228.5 & 229.9 & $+0.6\%$ \\
\midrule
\multicolumn{4}{l}{\emph{Architecture substitution (\texttt{small}$\rightarrow$\texttt{tiny}): $\approx -310$ GPU-h vs.\ the original plan.}} \\
\bottomrule
\end{tabular}
\end{table}

\subsection{Limitations}
The conclusions are bounded by the experimental scope. The five-axis pipeline is
a synthetic mathematical model and omits wavelength-dependent diffraction,
mixing-element nonlinearity, and detector dead-pixels of real THz hardware;
CIFAR-10 and MNIST are small benchmarks, so the per-axis attribution is not
validated at ImageNet scale or on true THz object categories; no learned
super-resolution baseline (a scoped-out Phase~E) was run, so the
``restore-resolution-first'' recommendation is not yet compared against a real
restoration pipeline; and the acceptance criteria are model-side, without a
human-perception upper bound for the extreme operating points.

\subsection{Recommendations for continuation}
\begin{enumerate}
  \item \textbf{Acquire real THz data.}\quad Even a small real dataset at the
        target frequency would let the next iteration calibrate the synthetic
        model against ground truth and renormalize the per-axis attribution.
  \item \textbf{Add a super-resolution pre-stage.}\quad Given that resolution
        dominates the attribution, the next ablation should place a learned
        upsampler in front of each backbone and re-run Phases~B2 / C2.
  \item \textbf{Enable graph compilation for TransNeXt.}\quad Rewriting the
        foveal-attention path as a compiled custom op would likely halve its
        inference latency and bring it inside the per-image budget.
  \item \textbf{Reproduce the attribution at scale.}\quad Repeat the per-axis
        analysis on a large robustness benchmark such as ImageNet-A
        [6] to test whether resolution dominance
        holds at full ImageNet scale.
  \item \textbf{Calibration intervention.}\quad The expected-calibration-error
        growth toward the severe levels indicates that a post-hoc temperature
        scaling step would improve the operational usefulness of the
        predictions without affecting top-1 accuracy.
\end{enumerate}
